\documentclass[../DoAn.tex]{subfiles}
\begin{document}


\section{Đặt vấn đề}
\label{section:1.1}
Sự phát triển mạnh mẽ của Internet vạn vật trong y tế (IoMT) đang tạo ra những bước tiến mang tính cách mạng trong công tác chăm sóc sức khỏe và theo dõi bệnh nhân. Tại các môi trường yêu cầu sự giám sát liên tục như viện dưỡng lão và bệnh viện, việc đảm bảo an toàn cho người cao tuổi và bệnh nhân gặp khó khăn về vận động luôn là một thách thức lớn đối với đội ngũ y tế. Đặc biệt, các sự cố té ngã xảy ra ngoài tầm quan sát của nhân viên nếu không được phát hiện và xử lý kịp thời có thể dẫn đến những hậu quả y khoa nghiêm trọng, thậm chí đe dọa đến tính mạng. Tuy nhiên, việc bố trí nhân sự túc trực 24/7 để theo dõi từng bệnh nhân là điều không khả thi do áp lực về nguồn nhân lực y tế ngày càng gia tăng. Thực trạng này đòi hỏi một giải pháp công nghệ tự động hóa có khả năng giám sát liên tục, phản hồi tức thời và quản lý tập trung, nhằm giảm thiểu tối đa rủi ro từ các tai nạn không mong muốn.

\section{Mục tiêu và phạm vi đề tài}
\label{section:1.2}
Hiện nay, trên thị trường đã xuất hiện một số giải pháp phát hiện té ngã, từ các hệ thống camera giám sát (Computer Vision) đến các thiết bị đeo tay thông minh (Smartwatch). Tuy nhiên, hệ thống camera thường xâm phạm quyền riêng tư nghiêm trọng và bị giới hạn bởi các điểm mù trong không gian, trong khi các thiết bị đeo cá nhân hiện hành phần lớn phụ thuộc vào điện thoại thông minh qua sóng Bluetooth, khiến chúng thiếu tính đồng bộ và khó quản lý tập trung trong môi trường viện dưỡng lão.

Dựa trên những hạn chế đó, đề tài hướng tới mục tiêu thiết kế và xây dựng một hệ thống cảnh báo khép kín chuyên dụng cho các cơ sở y tế tập trung. Cụ thể, hệ thống bao gồm thiết bị đeo thu thập dữ liệu trực tiếp từ người dùng và một nền tảng Web quản trị. Phạm vi của đồ án tập trung vào việc nghiên cứu và ứng dụng trí tuệ nhân tạo để phân tích trực tiếp hành vi trên thiết bị nhằm phát hiện các biến cố té ngã, đồng thời đảm bảo khả năng truyền nhận dữ liệu độc lập và liên tục mà không cần thông qua điện thoại thông minh của bệnh nhân.

\section{Định hướng giải pháp}
\label{section:1.3}
Để giải quyết bài toán trên, đồ án đề xuất giải pháp "Nghiên cứu và xây dựng hệ thống giám sát hành vi người già và phát hiện té ngã ứng dụng IoT và TinyML". 
Về mặt phần cứng thiết bị đeo, đồ án sử dụng kiến trúc học máy tại biên (Edge AI) thông qua nền tảng TinyML trên vi điều khiển ESP32-S3, cho phép phân tích dữ liệu cảm biến gia tốc và nhận diện hành vi ngay tại thiết bị với độ trễ siêu thấp. Về mặt truyền thông mạng, thiết bị sử dụng mạng di động 4G LTE kết hợp giao thức MQTT để đẩy dữ liệu trực tiếp lên máy chủ, loại bỏ sự phụ thuộc vào hạ tầng mạng Wi-Fi hay thiết bị trung gian. Cuối cùng, một hạ tầng Web Dashboard thời gian thực được xây dựng để cung cấp giao diện quản lý tập trung danh sách thiết bị và tự động phát âm thanh cảnh báo tức thời cho nhân viên chăm sóc.

Đóng góp chính của đồ án là việc hoàn thiện một hệ thống giám sát toàn diện từ phần cứng nhúng đến ứng dụng Web, giúp giảm tải áp lực cho nhân viên y tế, nâng cao hiệu quả phản hồi khẩn cấp và bảo vệ an toàn cho người cao tuổi.

\section{Bố cục đồ án}
\label{section:1.4}
Phần còn lại của báo cáo đồ án tốt nghiệp này được tổ chức như sau.

Chương 2 trình bày chi tiết về quá trình khảo sát các nghiên cứu và sản phẩm tương tự hiện có trên thị trường, đồng thời đưa ra các đặc tả chức năng cốt lõi của hệ thống phần mềm.

Chương 3 đi sâu vào việc giới thiệu các nền tảng lý thuyết và công nghệ được áp dụng trong đồ án, bao gồm kiến trúc truyền thông IoT (MQTT qua 4G LTE), nền tảng học máy TinyML, thuật toán mạng nơ-ron tích chập thặng dư 1 chiều (1D ResNet) và các công nghệ phát triển ứng dụng Web.

Chương 4 trình bày quá trình phân tích, thiết kế và xây dựng hệ thống, bao gồm kiến trúc phần mềm firmware theo mô hình Hướng sự kiện và Máy trạng thái, thiết kế cơ sở dữ liệu kép, giao thức truyền thông, cùng quá trình triển khai các thành phần cốt lõi như bộ lọc Kalman, pipeline TinyML, Backend và Frontend. Chương kết thúc bằng kết quả kiểm thử tích hợp hệ thống đầu cuối.

Chương 5 trình bày chi tiết ba đóng góp kỹ thuật nổi bật của đồ án: (i)~chiến lược thiết kế nhãn phân loại 5 lớp nhằm giảm thiểu cảnh báo sai từ các hành vi chuyển tư thế thường ngày; (ii)~tiến trình tối ưu hóa kiến trúc mô hình học sâu từ CNN-LSTM đến mạng 1D ResNet tương thích TinyML với các kỹ thuật tối ưu phần cứng ESP-NN; và (iii)~phương pháp lấy mẫu IMU chính xác sử dụng ngắt phần cứng PCNT thay thế cơ chế trễ phần mềm FreeRTOS.

Cuối cùng, Chương 6 tổng kết lại các kết quả chính đã đạt được, đồng thời đề xuất những hướng phát triển và mở rộng ứng dụng trong tương lai.

\end{document}